\documentclass[11pt]{article}
% Enable LaTeX->HTML conversion
\usepackage[mathjax]{lwarp}
\usepackage[margin=1in]{geometry}
\usepackage{amsmath,amsfonts,amssymb}

\title{Dirac Notes for Integration over Energy Surfaces}
\author{}
\date{\today}

\begin{document}
\maketitle
\tableofcontents

\section{General Concept}
The Dirac delta distribution $\delta(x)$ is defined by its action on a sufficiently smooth test function $f(x)$ under an integral:
\[
\int_{-\infty}^{\infty} f(x) \delta(x) \, dx = f(0).
\]

\section{Shifting Properties}
The Dirac delta can be seen as an operator that "picks out" the value of a function at a specific point. For a function \( f(x) \) and a point \( a \), we consider:
\[
\int_{-\infty}^{\infty} f(x) \delta(x - a) \, dx.
\]
By substituting \( u = x - a \), we have \( dx = du \). As \( x \to \pm\infty \), \( u \) also goes to \( \pm\infty \). The integral becomes:
\[
\int_{-\infty}^{\infty} f(u + a) \delta(u) \, du = f(0 + a) = f(a).
\]
This result is fundamental and is often referred to as the \textbf{sifting} or \textbf{sampling property}.

\section{Composition with Functions}
\subsection{Linear Scaling}
Consider a linear scaling of the argument $ax$ within the delta function:
\[
\int_{-\infty}^{\infty} f(x) \delta(ax) \, dx.
\]
Substituting \( u = a x \), which implies \( dx = \frac{du}{a} \), we have:
\[
\int_{-\infty}^{\infty} f(x) \delta(a x) \, dx = \int_{-\infty}^{\infty} f\left(\frac{u}{a}\right) \delta(u) \frac{du}{a}.
\]
Since the Dirac delta only evaluates the function at \( u = 0 \), it fulfills the symmetry property \( \delta(u) = \delta(-u) \). This ensures that the result depends only on the magnitude of the scaling factor, yielding:
\[
\int_{-\infty}^{\infty} f(x) \delta(a x) \, dx = \frac{1}{|a|} f(0).
\]
\subsection{Composition with a General Function}
If the argument is a non-linear function $g(x)$, the delta function triggers only at the roots $x_i$ where $g(x_i) = 0$. 
Assuming first a single root $x_0$, we linearize $g(x)$ around this point using a Taylor expansion: $g(x) \approx g'(x_0)(x - x_0)$. Using the scaling rule from above, we find:
\[
\int_{-\infty}^{\infty} f(x) \delta(g(x)) \, dx \approx \int_{-\infty}^{\infty} f(x) \delta(g'(x_0)(x - x_0)) \, dx = \frac{1}{|g'(x_0)|} f(x_0).
\]
If $g(x)$ has multiple isolated roots $x_i$, the integral generalizes to a sum over all these points:
\[
\int_{-\infty}^{\infty} f(x) \delta(g(x)) \, dx = \sum_i \frac{f(x_i)}{|g'(x_i)|}.
\]      
\subsection{Integration over Energy Surfaces}
The dispersion relation $E(\mathbf{k})$ defines energy surfaces in momentum space. To integrate a function $f(\mathbf{k})$ over a constant energy surface $E_0$, we use the delta function:
\[
\int f(\mathbf{k}) \delta(E(\mathbf{k}) - E_0) \, d\mathbf{k}.
\]
The delta function restricts the integration to points where $E(\mathbf{k}) = E_0$. 
The increase of the energy is given locally by its gradient $\nabla_{\mathbf{k}} E(\mathbf{k})$. A Taylor expansion around a point $\mathbf{k}_0$ on the energy surface yields:
\[
E(\mathbf{k}) \approx E(\mathbf{k}_0) + \nabla_{\mathbf{k}} E(\mathbf{k}_0) \cdot (\mathbf{k} - \mathbf{k}_0).
\]
Only the component of $\mathbf{k}$ parallel to the gradient (normal to the surface) contributes to the shift in energy. Applying the scaling property from the 1D case to this normal direction, we find:
\[
\int f(\mathbf{k}) \delta(E(\mathbf{k}) - E_0) \, d\mathbf{k} = \int_{E(\mathbf{k}) = E_0} \frac{f(\mathbf{k})}{|\nabla_{\mathbf{k}} E(\mathbf{k})|} \, dS,
\]
where $dS$ is the surface element on the energy surface defined by $E(\mathbf{k}) = E_0$.








\end{document}
